% TEMPLATE-MILC-v3.tex - plantilla maestra UNICA de las guias MILC v3.
%
% La usan las tres familias de guias, cada una con su driver:
%   · Grados 6-11  -> scripts/build-guias-g11.py
%   · Bebras       -> scripts/build-guias-bebras.py
%   · Semillero    -> scripts/build-guias-semillero.py
%
% Antes habia tres copias casi identicas (~400 lineas iguales, 6 distintas):
% cada mejora habia que aplicarla tres veces, y en la practica no se hacia
% (el motor de recursos vivia solo en dos de ellas). Lo que varia entre
% familias esta parametrizado en 6 placeholders que cada driver llena
% (se nombran aqui SIN su sintaxis real: el builder reemplaza por texto plano,
% tambien dentro de los comentarios, y un valor multilinea rompe el preambulo):
%
%   HEADER_IZQ        encabezado izquierdo de pagina
%   HEADER_DER        encabezado derecho de pagina
%   PORTADA_ETIQUETA  rotulo sobre el numero grande (GRADO/MOMENTO/MODULO)
%   PORTADA_SERIE     linea de serie de la portada
%   PORTADA_ID        identificador de la guia en la portada
%   PORTADA_CIRCULO   el circulito de grado (°), o vacio si no aplica
%
% El resto de claves las documenta content/guias/_SCHEMA.md. El script valida
% que no queden placeholders sin reemplazar antes de escribir el archivo final.

\documentclass[11pt,letterpaper]{article}
\usepackage[margin=2.0cm]{geometry}
\usepackage[table]{xcolor}
\usepackage{fontspec}
\usepackage[spanish]{babel}
\usepackage{tikz}
% Librerías TikZ para los diagramas de las guías (bloque `recursos:`):
%  · babel  → babel en español vuelve activos caracteres como «>», y TikZ los usa
%             en sus opciones (>=stealth, ->). Sin esta librería, cualquier
%             diagrama con flechas revienta con «\language@active@arg> has an
%             extra }». Debe ir ANTES de que se dibuje nada.
%  · positioning        → sintaxis «below left=2pt and 3pt».
%  · decorations.pathreplacing → llaves ({brace}) para señalar rangos.
\usetikzlibrary{babel,positioning,decorations.pathreplacing}
\usepackage{graphicx}
% Circuitos: el trabajo de electrónica se hace hoy con micro:bit y sus sensores
% integrados (MakeCode, sin cableado), así que no se necesitan esquemáticos.
% Si algún día entran montajes con protoboard, basta descomentar esta línea:
% los diagramas .tex/.tikz declarados en `recursos:` ya se incrustan solos.
% \usepackage{circuitikz}
\usepackage[most]{tcolorbox}
\usepackage{tabularx}
\usepackage{array}
\usepackage{fancyhdr}
\usepackage{titlesec}
\usepackage[document]{ragged2e}  % bandera a la derecha en todo el cuerpo
\usepackage{enumitem}
\usepackage{microtype}

% Helvetica Neue no trae la flecha U+2192, asi que cada «→» escrito literalmente
% en el YAML se compilaba a NADA, en silencio (462 flechas en 61 guias: rutas del
% tipo «paleta → tipografia → lockup» salian rotas). Mapeamos el caracter a su
% version matematica, que si tiene glifo. Asi el autor puede seguir escribiendo
% «→» comodamente en el YAML.
\usepackage{newunicodechar}
\newunicodechar{→}{\ensuremath{\rightarrow}}
% Mismo problema con los simbolos de marca y vinetas: el «✓» de «una palomita ✓»
% salia como cuadrito vacio en el PDF de todas las guias que lo usaban.
\usepackage{pifont}
\newunicodechar{✓}{\ding{51}}
\newunicodechar{✗}{\ding{55}}
\newunicodechar{✔}{\ding{52}}
\newunicodechar{•}{\textbullet}
\newunicodechar{≥}{\ensuremath{\geq}}
\newunicodechar{≤}{\ensuremath{\leq}}

\setmainfont{Helvetica Neue}
\setsansfont{Helvetica Neue}
\setlength{\parindent}{0pt}
\setlength{\parskip}{6pt}
% Sin particion de palabras: en 6.º-7.º una palabra cortada por guion al final
% de linea («Comuni-cacion») frena la lectura mucho mas que una linea corta.
\hyphenpenalty=10000
\exhyphenpenalty=10000
\pretolerance=2000
\tolerance=3000
\setlength{\emergencystretch}{3em}
\linespread{1.10}
\renewcommand{\arraystretch}{1.25}
\setlist{leftmargin=5mm,itemsep=1mm,topsep=1mm}
\newcolumntype{Y}{>{\RaggedRight\arraybackslash}X}

\definecolor{milcMagenta}{HTML}{E6007E}
\definecolor{milcVino}{HTML}{5A0038}
\definecolor{milcTurquesa}{HTML}{0093A5}
\definecolor{milcVerde}{HTML}{58A923}
\definecolor{milcMostaza}{HTML}{E5B400}
\definecolor{milcOcre}{HTML}{B86600}
\definecolor{milcNegro}{HTML}{241816}
\definecolor{milcGris}{HTML}{F7F7F4}
\definecolor{milcLinea}{HTML}{E8E2DD}
\definecolor{milcGradePrimary}{HTML}{0D9488}
\definecolor{milcGradeDark}{HTML}{115E59}
\definecolor{milcGradeSoft}{HTML}{CCFBF1}
\definecolor{milcGradeAccent}{HTML}{CFE7FF}
\definecolor{milcGradeText}{HTML}{FFFFFF}
\color{milcNegro}

\pagestyle{fancy}
\fancyhf{}
\lhead{\small Territorio interior · Educación socioemocional}
\rhead{\small Grado 8° · Momento 4 · MILC v3}
\cfoot{\small\thepage}
\renewcommand{\headrulewidth}{0pt}

\newtcolorbox{titlebox}[1]{
  enhanced,colback=#1,colframe=#1,arc=7mm,boxrule=0pt,
  left=7mm,right=7mm,top=3mm,bottom=3mm,
  before skip=8pt,after skip=8pt,
  fontupper=\sffamily\bfseries\Large\color{white}
}

\newtcolorbox{softbox}[3]{
  enhanced,colback=#2,colframe=#1,borderline west={4pt}{0pt}{#1},
  arc=2mm,boxrule=.4pt,left=4mm,right=4mm,top=3mm,bottom=3mm,
  before skip=6pt,after skip=8pt,title={#3},coltitle=white,
  colbacktitle=#1,fonttitle=\sffamily\bfseries,fontupper=\RaggedRight,
  attach boxed title to top left={xshift=3mm,yshift=-2mm},
  boxed title style={arc=2mm, boxrule=0pt}
}

\newtcolorbox{drawbox}[2]{
  enhanced,colback=white,colframe=#1,arc=4mm,boxrule=1pt,
  left=5mm,right=5mm,top=4mm,bottom=4mm,before skip=8pt,after skip=8pt,
  title={#2},coltitle=white,colbacktitle=#1,fonttitle=\sffamily\bfseries,
  fontupper=\RaggedRight,
  attach boxed title to top left={xshift=4mm,yshift=-2mm},
  boxed title style={arc=3mm, boxrule=0pt}
}

\newtcolorbox{infoband}[1]{
  enhanced,colback=#1!8,colframe=#1!50,arc=2mm,boxrule=.5pt,
  left=4mm,right=4mm,top=2mm,bottom=2mm,before skip=4pt,after skip=4pt,
  fontupper=\small\RaggedRight
}

% Puente narrativo entre secciones (contrato editorial Conéctate).
% Da continuidad: "Acabas de X. Ahora vas a Y porque Z."
% Visualmente: banda izquierda en color del grado, texto en itálica pequeña.
\newtcolorbox{puentebox}{
  enhanced,colback=milcGradeSoft,colframe=milcGradeDark,
  borderline west={3pt}{0pt}{milcGradeDark},
  arc=0mm,boxrule=0pt,left=5mm,right=4mm,top=2.5mm,bottom=2.5mm,
  before skip=4pt,after skip=8pt,
  fontupper=\itshape\small\color{milcNegro}\RaggedRight
}
% La flecha va en modo matematico a proposito: Helvetica Neue NO tiene el glifo
% U+2192, asi que \textrightarrow se compilaba a NADA («Missing character: There
% is no → in font Helvetica Neue») y el puente salia sin flecha. En modo
% matematico la flecha viene de la fuente matematica y si se dibuja.
\newcommand{\puente}[1]{%
  \begin{puentebox}%
    \textcolor{milcGradeDark}{$\rightarrow$}\hspace{2mm}#1%
  \end{puentebox}%
}

\newcommand{\linead}[1]{\noindent\color{milcLinea}\rule{#1}{0.5pt}\color{milcNegro}}
\newcommand{\respuesta}[1]{\par\vspace{2mm}\foreach \n in {1,...,#1}{\noindent\color{milcLinea}\rule{\linewidth}{0.5pt}\par\vspace{5mm}}\color{milcNegro}}
\newcommand{\checkbox}{\textcolor{milcGradeDark}{\fbox{\phantom{x}}}\hspace{5pt}}

% ─── Listas aireadas para las guías socioemocionales ────────────────────
% Componer las verificaciones como «\checkbox texto\\» deja el texto que salta
% de línea debajo del cuadrito y apelmaza el bloque. Estos entornos alinean la
% continuación y airean los ítems. Los usa Territorio Interior; cualquier
% familia puede hacerlo.
\newlist{checklista}{itemize}{1}
\setlist[checklista]{
  label=\textcolor{milcGradeDark}{\fbox{\phantom{x}}},
  leftmargin=9mm, labelsep=3.2mm, itemsep=2.4mm,
  topsep=2.5mm, parsep=0pt, partopsep=0pt, before=\RaggedRight
}
\newlist{pasoslista}{enumerate}{1}
\setlist[pasoslista]{
  label=\textcolor{milcGradeDark}{\sffamily\bfseries\arabic*.},
  leftmargin=8mm, labelsep=3mm, itemsep=2.8mm,
  topsep=1mm, parsep=0pt, partopsep=0pt, before=\RaggedRight
}

\newcommand{\phasecard}[4]{
\begin{tcolorbox}[enhanced,colback=#3,colframe=#2,arc=3mm,boxrule=.5pt,height=3.05cm,valign=center,halign=center,left=1.5mm,right=1.5mm,top=2mm,bottom=2mm]
{\sffamily\bfseries\fontsize{22}{24}\selectfont\textcolor{#2}{#1}}\par
{\sffamily\bfseries\small #4}
\end{tcolorbox}
}

% ── Piezas del rediseno 2026 (legibilidad 6.º-11.º) ───────────────────
% Banda de estacion: reemplaza el titlebox macizo. Titulo a la izquierda,
% dato practico (tiempo/modalidad) a la derecha, en una sola linea.
\newcommand{\milcestacion}[2]{%
\begin{tcolorbox}[enhanced,colback=#1,colframe=#1,arc=2mm,boxrule=0pt,
  left=5mm,right=5mm,top=2.6mm,bottom=2.6mm,before skip=11pt,after skip=7pt]
{\sffamily\bfseries\fontsize{15}{18}\selectfont\color{white}#2}
\end{tcolorbox}}

% Tarjeta de paso numerada: sustituye las tablas «Pilar / Aplicacion», que
% eran formato de consulta docente, no de estudiante. El numero va en un
% circulo sobre el borde para que la secuencia se lea de un vistazo.
\newcommand{\milcpaso}[3]{%
\begin{tcolorbox}[enhanced,colback=white,colframe=milcLinea,arc=1.5mm,boxrule=.9pt,
  left=14mm,right=4mm,top=3mm,bottom=3mm,before skip=4pt,after skip=4pt,
  fontupper=\RaggedRight,
  overlay={\node[anchor=north west,circle,fill=#2,text=white,inner sep=0pt,
    minimum size=7.4mm,font=\sffamily\bfseries\small]
    at ([xshift=3.6mm,yshift=-3.2mm]frame.north west) {#1};}]
#3
\end{tcolorbox}}

% Casilla marcable de tamano util con lapiz (la anterior era \fbox{\phantom{x}},
% de alto variable segun el texto que la seguia).
\newcommand{\milcmarca}{\framebox[3.8mm]{\rule{0pt}{3.4mm}}}

% Fila marcable de lista suelta (checks de la estacion 1).
\newcommand{\milccheckrow}[1]{%
\par\vspace{1.8mm}\noindent
\makebox[6.4mm][l]{\raisebox{-.55mm}{\textcolor{milcNegro}{\milcmarca}}}%
\parbox[t]{\dimexpr\linewidth-6.4mm\relax}{\RaggedRight #1}\par}

% Celda marcable dentro de la rubrica (tabla de autoevaluacion). Misma
% sangria francesa, pero pensada para vivir dentro de una columna X.
\newcommand{\milcopcion}[1]{%
\makebox[5.2mm][l]{\raisebox{-.55mm}{\textcolor{milcNegro}{\milcmarca}}}%
\parbox[t]{\dimexpr\linewidth-5.2mm\relax}{\RaggedRight #1}}

% ── Recursos visuales de la guía ──────────────────────────────────────
% Declarados en el bloque `recursos:` del YAML y emitidos por el builder.
% \guiaFigura   → imagen rasterizada o PDF (foto, carta celeste, montaje).
% \guiaDiagrama → fuente TikZ/circuitikz (.tex) incrustada como vector:
%                 es la vía para circuitos y esquemáticos editables.
% #1 = ruta relativa al .tex · #2 = pie (puede ir vacío)
\newcommand{\guiaFigura}[2]{%
\begin{center}
\includegraphics[width=0.88\linewidth,height=0.38\textheight,keepaspectratio]{#1}\\[1.5mm]
{\footnotesize\itshape\textcolor{milcNegro!75}{#2}}
\end{center}
}

\newcommand{\guiaDiagrama}[2]{%
\begin{center}
\input{#1}\\[1.5mm]
{\footnotesize\itshape\textcolor{milcNegro!75}{#2}}
\end{center}
}

\begin{document}
\thispagestyle{empty}

% ── Portada ───────────────────────────────────────────────────────────
% Antes: un rectangulo de color a pagina completa. Se veia bien en pantalla,
% pero en la fotocopiadora del colegio se comia el toner y salia gris sucio.
% Ahora la banda ocupa el tercio superior y el resto va en blanco.
\begin{tikzpicture}[remember picture,overlay]
  \path[fill=milcGradePrimary]
    (current page.north west) rectangle ([yshift=-10.6cm]current page.north east);

  \node[anchor=north west,text=milcGradeText] at ([xshift=2.0cm,yshift=-1.55cm]current page.north west)
    {\sffamily\bfseries\fontsize{12}{14}\selectfont MOMENTO};
  \node[anchor=north west,text=milcGradeText,opacity=.85] at ([xshift=2.0cm,yshift=-2.35cm]current page.north west)
    {\sffamily\bfseries\fontsize{8.5}{10}\selectfont TERRITORIO INTERIOR · SOCIOEMOCIONAL};
  \node[anchor=north east,text=milcGradeText,opacity=.85,align=right] at ([xshift=-2.0cm,yshift=-1.55cm]current page.north east)
    {\sffamily\bfseries\fontsize{9}{12}\selectfont I.E. Sor María Juliana\\Cartago, Valle del Cauca};

  \node[anchor=west,text=milcGradeText] (gradonum) at ([xshift=1.85cm,yshift=-7.1cm]current page.north west)
    {\sffamily\bfseries\fontsize{112}{112}\selectfont 4};
  % El circulito de grado (°) solo aplica a las guias de grado («11°»); en
  % Bebras y Semillero el numero grande es un momento/modulo, no un grado.
  % Se posiciona relativo al nodo `gradonum`, no en coordenadas absolutas.
  

  \node[anchor=west,text=milcGradeText,text width=11.4cm,align=left] at ([xshift=1.5cm,yshift=0.5cm]gradonum.east)
    {
     {\sffamily\bfseries\fontsize{9}{11}\selectfont\MakeUppercase{Territorio interior · Socioemocional}}\\[3mm]
     {\sffamily\bfseries\fontsize{21}{25}\selectfont\setlength{\baselineskip}{27pt}Rumor y\\[4pt]reputación\\[4pt]deshacer cuesta más que hacer}\\[3.5mm]
     {\sffamily\bfseries\fontsize{15}{18}\selectfont Grado 8° · Momento 4}
    };
\end{tikzpicture}

\vspace*{9.4cm}
\vfill

{\sffamily\bfseries\fontsize{8}{10}\selectfont\textcolor{milcGradeDark}{LO QUE TE LLEVAS HOY}}

\begin{tcolorbox}[enhanced,colback=white,colframe=milcNegro,arc=2mm,boxrule=1.6pt,
  left=5mm,right=5mm,top=4mm,bottom=4mm,before skip=3pt,after skip=12pt,
  fontupper=\RaggedRight\fontsize{11}{15.5}\selectfont]
Tu cadena del rumor: un rumor real reconstruido hacia atrás hasta donde alcance —quién te lo dijo, quién se lo dijo—, la cuenta de cuántos lo saben frente a cuántos saben la corrección, y tu regla de tres preguntas antes de repetir algo sobre alguien.
\end{tcolorbox}

\begin{tcolorbox}[enhanced,colback=milcGris,colframe=milcGris,arc=2mm,boxrule=0pt,
  left=5mm,right=5mm,top=3.5mm,bottom=3.5mm,after skip=12pt]
{\sffamily\bfseries\fontsize{8}{10}\selectfont\textcolor{milcNegro!55}{ESTA GUÍA ES DE}}\\[3mm]
\begin{tabularx}{\linewidth}{X p{3.2cm} p{3.2cm}}
\linead{\linewidth} & \linead{3.0cm} & \linead{3.0cm}\\[-1mm]
{\fontsize{8.5}{10}\selectfont\textcolor{milcNegro!55}{Nombre completo}} &
{\fontsize{8.5}{10}\selectfont\textcolor{milcNegro!55}{Grupo}} &
{\fontsize{8.5}{10}\selectfont\textcolor{milcNegro!55}{Fecha}}\\
\end{tabularx}
\end{tcolorbox}

\vfill

\noindent\textcolor{milcLinea}{\rule{\linewidth}{1pt}}\\[2mm]
{\sffamily\bfseries\fontsize{9.5}{12}\selectfont Autor: Dr. Álvaro Cárdenas Orozco}\\[1mm]
{\sffamily\fontsize{8.5}{11}\selectfont\textcolor{milcNegro!55}{Metodología MILC · Apertura ancestral · Rumor y reputación · Triángulo Dussel-Estoicismo-Floridi}}

\newpage

\section*{Apertura: Rumor y reputación: lo que cuesta una fama y lo que cuesta deshacerla}

\begin{softbox}{milcGradeDark}{milcGradeSoft}{Saber ancestral}
En los talleres de calado de Cartago hay una destreza que no se luce y que las caladoras valoran tanto como bordar: \textbf{saber remendar}. La antropóloga \textbf{Tania Pérez-Bustos} lo documentó trabajando con ellas, y recogió el criterio en una frase de Olivia, caladora: \emph{«no quieres que se vea remendada»} (Pérez-Bustos, 2019). \textbf{Ahí está todo el oficio del arreglo:} un remiendo bueno \textbf{es el que no se nota}; si se nota, la pieza queda marcada para siempre, y quien la mire va a mirar primero el remiendo y después el bordado. \textbf{Y hay algo más, que es lo duro:} las caladoras saben reconocer \textbf{cuándo el daño ya es demasiado grande}. Cuando la tela cedió más de la cuenta, \textbf{no hay remiendo posible} ---y decirlo a tiempo también es parte del oficio---. \emph{Ahora traduce.} \textbf{Una reputación funciona exactamente así:} se sostiene en lo que \textbf{no se nota}, un arreglo mal hecho la marca más que el daño original, \textbf{y tiene un punto sin retorno}. \emph{La diferencia es que la tela la dañó alguien con las manos; una fama la dañan veinte personas repitiendo, y ninguna cree haber hecho nada.}
\end{softbox}

\begin{softbox}{milcTurquesa}{milcGris}{Saber contemporáneo}
¿Por qué es tan difícil quitarse una fama? \textbf{Porque el cerebro no trata igual lo bueno y lo malo.} \textbf{Roy Baumeister} y su equipo revisaron centenares de estudios y llegaron a una conclusión que resumieron en el título: \textbf{lo malo es más fuerte que lo bueno}. Los hechos malos pesan más que los buenos en la vida diaria, en las relaciones y en el aprendizaje; \textbf{la información mala se procesa más a fondo}; y ---esto es lo que nos importa hoy--- \textbf{las malas impresiones y los malos estereotipos se forman más rápido y son más resistentes a que los desmientan} que los buenos (Baumeister, Bratslavsky, Finkenauer y Vohs, 2001). \textbf{Traducido a números que se sienten:} \emph{basta una cosa para que te pongan una fama, y no bastan diez para quitártela}. \textbf{Y hay un segundo mecanismo que lo empeora:} el rumor \textbf{viaja} porque es interesante; \textbf{la corrección no viaja} porque es aburrida. \emph{Nadie reenvía «resulta que era mentira».} \textbf{Por eso el número de gente que sabe la versión falsa siempre es mayor que el de la que sabe la verdadera} ---y esa diferencia, sin que nadie haga nada, no se cierra sola---.
\end{softbox}

\begin{softbox}{milcMostaza}{milcGris}{Pregunta-puente}
\textit{Una caladora sabe \textbf{cuándo un daño ya no tiene remiendo}, y lo dice a tiempo. \textbf{¿De qué persona de tu colegio sabes algo que nunca comprobaste} \ldots \textbf{y cuántas veces lo has repetido}?}
\end{softbox}

\begin{softbox}{milcVerde}{milcGris}{Saber y saber hacer}
Al terminar podrás: (1) \textbf{explicar} por qué una mala impresión se forma más rápido y aguanta más que una buena; (2) \textbf{reconstruir} una cadena de rumor hacia atrás y descubrir que casi nunca llega a una fuente; (3) \textbf{comparar} cuánta gente sabe la versión falsa frente a cuánta sabe la corrección, y entender por qué esa brecha no se cierra sola; (4) \textbf{usar} una regla de tres preguntas antes de repetir algo sobre alguien, y saber cómo se corrige \textbf{donde} se dijo.
\end{softbox}



\small
\begin{tabularx}{\textwidth}{>{\bfseries}p{3.0cm}Y}
\rowcolor{milcGradePrimary!12}Autor & Dr. Álvaro Cárdenas Orozco\\
DBA & Comprende los fenómenos que afectan a su territorio, regula sus emociones ante ellos y acompaña a otros con empatía (transversal 6°--11°, articulado con la política «Escuela, territorio de vida» del MEN, Resolución 006519 de 2025, y la Ley 1620 de 2013)\\
Referentes & Primera ayuda psicológica (OMS, 2011/2012) · Directrices IASC (2007) · USGS y Servicio Geológico Colombiano · Pueblo nasa y la minga (CRIC) · Dussel · Estoicismo · Floridi\\
Producto final & Tu cadena del rumor: un rumor real reconstruido hacia atrás hasta donde alcance —quién te lo dijo, quién se lo dijo—, la cuenta de cuántos lo saben frente a cuántos saben la corrección, y tu regla de tres preguntas antes de repetir algo sobre alguien.\\
\end{tabularx}
\normalsize

\puente{En el momento 3 aprendiste qué hacer cuando ves algo. \textbf{Hoy el daño no se ve:} se hace hablando, y el que lo hace casi nunca siente que hizo nada.}

\milcestacion{milcGradeDark}{Ruta MILC: así vamos a aprender}
\begin{tabularx}{\textwidth}{>{\centering\arraybackslash}X>{\centering\arraybackslash}X>{\centering\arraybackslash}X>{\centering\arraybackslash}X}
\phasecard{1}{milcVerde}{milcGris}{Escucha\\[-1mm]\small Observo, escucho y pregunto.} &
\phasecard{2}{milcTurquesa}{milcGris}{Sistematización\\[-1mm]\small Rastreo y no repito.} &
\phasecard{3}{milcMagenta}{milcGris}{Praxis\\[-1mm]\small Construyo y aplico.} &
\phasecard{4}{milcVino}{milcGris}{Evaluación\\[-1mm]\small Reflexión liberadora.}
\end{tabularx}


\puente{Primero, un rumor real ---uno que circule ahora--- y la pregunta que casi nadie se hace: \emph{¿quién me lo dijo a mí?}}

\milcestacion{milcVerde}{Estación 1 de 3 · Escucha}

\begin{softbox}{milcVerde}{milcGris}{Escena para escuchar}
\textbf{Actividad 1 · IDENTIFICA --- La cadena.} \textbf{Sin nombres propios: usa iniciales.} \textbf{Primera parte.} Piensa en \textbf{algo que «se dice» de alguien} en tu colegio ---no lo peor: sirve cualquier cosa, incluso algo que parece inofensivo---. Escríbelo tal como circula. \textbf{Segunda parte, la del ejercicio.} Reconstruye la cadena \textbf{hacia atrás}: \emph{¿quién te lo dijo a ti? ¿Y a esa persona quién se lo dijo? ¿Y a esa?} \textbf{Llega tan atrás como puedas} y anota \textbf{dónde se te acabó la cadena}. \textbf{Tercera parte.} ¿Alguna vez alguien \textbf{comprobó} eso? ¿Se lo preguntaron a la persona de la que se trata? \textbf{Y la incómoda:} ¿tú se lo has repetido a alguien? ¿A cuántos? \emph{Cuenta, aunque dé pena: la cuenta es el punto de toda la actividad.}
\end{softbox}

{\sffamily\bfseries\fontsize{8}{10}\selectfont\textcolor{milcNegro!55}{MARCA CUANDO LO HAYAS HECHO}}
\milccheckrow{Escribiste el rumor tal como circula, sin corregirlo.}
\milccheckrow{Reconstruiste la cadena hacia atrás y anotaste \textbf{dónde se acabó}.}
\milccheckrow{Contaste a cuántas personas se lo has repetido tú.}

\begin{infoband}{milcVerde}


\textbf{En tu cuaderno (Actividad 1 · IDENTIFICA).} Título: «Actividad 1. La cadena». \textbf{Formato:} el rumor + la cadena hacia atrás con iniciales + dónde se acabó + a cuántos se lo repetiste tú. \textbf{Extensión:} media página. \textbf{Sabes que terminaste cuando} puedes decir \textbf{en qué eslabón se te perdió el rastro}. Esta página es tuya: \textbf{nadie más tiene que leerla si tú no quieres}.
\end{infoband}

\puente{Ya lo tienes. Ahora, por qué lo malo pesa más, por qué la corrección no viaja y qué significa que un daño ya no tenga remiendo.}

\milcestacion{milcTurquesa}{Fase 2 · Sistematización: entender la fama}

\begin{softbox}{milcTurquesa}{milcGris}{Cuatro claves sobre el rumor}
Cuatro claves para entender por qué una fama se pone en un día y no se quita en un año.
\end{softbox}

\milcpaso{1}{milcTurquesa}{\textbf{Lo malo pesa más, y no es impresión tuya.} Baumeister y su equipo lo revisaron a fondo: los hechos malos pesan más que los buenos, \textbf{la información mala se procesa más a fondo}, y \textbf{las malas impresiones se forman más rápido y resisten más los desmentidos} que las buenas (Baumeister, Bratslavsky, Finkenauer y Vohs, 2001). \emph{Piénsalo con tu propia experiencia:} de alguien que te cayó mal una vez te acuerdas años; de alguien que te cayó bien, a veces ni el nombre. \textbf{Consecuencia práctica y dura:} \emph{una sola cosa alcanza para ponerle a alguien una fama, y diez cosas buenas no alcanzan para quitársela}. \textbf{Por eso el que repite un rumor no está «solo comentando»: está poniendo algo que pesa el doble y sale la mitad de fácil.}}
\milcpaso{2}{milcTurquesa}{\textbf{El rumor viaja; la corrección no.} Un rumor circula porque es \textbf{interesante} ---sorprende, indigna, da de qué hablar---. \emph{Una corrección es lo contrario: es aburrida.} \textbf{Nadie escribe en un grupo «acuérdense de lo que dijimos ayer: resulta que era mentira».} \textbf{Y de ahí sale el dato que hay que tener en la cabeza:} el número de gente que sabe la versión falsa \textbf{siempre} es mayor que el de la que sabe la corrección. \emph{Esa brecha no se cierra sola con el tiempo:} se cierra solo si alguien hace el trabajo aburrido de corregir \textbf{donde} se dijo, y con la misma insistencia. \textbf{Casi nadie lo hace, y por eso las famas duran.}}
\milcpaso{3}{milcTurquesa}{\textbf{Nadie inventó nada, y ese es el problema.} Cuando se rastrea una cadena hacia atrás casi nunca aparece un culpable: \emph{aparece que «me dijo fulano», «yo lo escuché», «todo el mundo lo sabe».} \textbf{La cadena se acaba en el aire.} \emph{Y ahí está la trampa:} como nadie lo inventó, \textbf{nadie se siente responsable} ---igual que en la difusión de responsabilidad del momento 3, el daño se repartió en pedacitos---. \textbf{Pero fíjate en la aritmética:} \emph{un rumor sin nadie que lo repita se muere en un día}. \textbf{No hace falta encontrar al que empezó: basta con que se acabe en ti.}}
\milcpaso{4}{milcTurquesa}{\textbf{Hay daños con remiendo y daños sin remiendo.} Las caladoras distinguen las dos cosas, y saber cuál es cuál es parte del oficio. \emph{Una fama recién puesta se puede desarmar} ---si alguien corrige rápido, donde se dijo y delante de los mismos---. \textbf{Pero cuando lleva meses, ya no es un rumor: es «lo que se sabe» de esa persona,} y para entonces ni siquiera hay a quién corregirle. \emph{Y esto es lo que hay que entender antes de repetir algo:} \textbf{no estás pasando una información, estás poniendo una puntada}. \textbf{Y de las puntadas, unas se descosen y otras no.}}

\begin{drawbox}{milcTurquesa}{Las tres preguntas antes de repetir algo sobre alguien}
\begin{pasoslista}
\item \textbf{¿Quién me lo dijo, y a esa persona quién?} \emph{Si la cadena se acaba en «todo el mundo lo sabe», no hay fuente.}
\item \textbf{¿Lo comprobó alguien con la persona de la que se trata?} \emph{Casi nunca. Y preguntar es gratis.}
\item \textbf{Si fuera falso, ¿podría deshacerse el daño?} \emph{Si la respuesta es no, ya sabes qué estás repitiendo.}
\item \textbf{Si alguna respuesta falla, no lo repitas.} No hay que discutir con nadie: \textbf{basta con que la cadena se acabe en ti}.
\end{pasoslista}
\end{drawbox}

\begin{softbox}{milcMostaza}{milcGris}{Errores comunes y su corrección}
\textbf{«Yo solo lo comenté»:} repetir es publicar; el que repite es el que lo lleva a gente nueva. \textbf{«Es que todo el mundo lo sabe»:} eso no es una fuente, es el final de una cadena rota. \textbf{«Yo no lo inventé»:} un rumor sin quien lo repita se muere en un día. \textbf{Corregir en privado un daño público:} la corrección tiene que ir donde fue el rumor. \textbf{Corregir una sola vez:} el rumor se dijo veinte veces; la corrección, una. \textbf{Preguntarle a todo el mundo menos a la persona:} preguntar es gratis. \textbf{Creer que el tiempo lo arregla:} lo malo resiste más que lo bueno; el tiempo solo lo asienta.
\end{softbox}

\begin{infoband}{milcTurquesa}


\textbf{En tu cuaderno (Actividad 2 · EXPLICA).} Título: «Actividad 2. Los dos números». \textbf{Formato:} para tu rumor, calcula \textbf{cuántas personas crees que saben la versión que circula} y \textbf{cuántas sabrían la corrección} si alguien la hiciera una vez. Escribe los dos números y explica la diferencia. \textbf{Extensión:} media página. \textbf{Sabes que terminaste cuando} puedes explicar por qué \textbf{una mala impresión aguanta más que una buena}.
\end{infoband}

\puente{Ya sabes cómo funciona. Es hora de cortar la cadena en el único punto donde puedes: \textbf{tú}.}

\milcestacion{milcMagenta}{Fase 3 · Praxis: cortar la cadena}

\begin{softbox}{milcMagenta}{milcGris}{Cómo se corta una cadena de la que uno hace parte}
Esto no se resuelve encontrando culpables. Se resuelve en un solo punto de la cadena, y es el único que controlas.
\end{softbox}

\milcpaso{1}{milcMagenta}{\textbf{La cadena hacia atrás, terminada (8 min).} Completa el rastreo hasta donde llegues y escribe en qué eslabón se acabó. \emph{Compara con dos compañeros:} \textbf{van a descubrir que casi todas las cadenas terminan igual}, en «me dijeron» o en «todo el mundo lo sabe». \textbf{Y a veces algo peor:} que la cadena \textbf{se cierra en círculo} ---A se lo oyó a B, B a C y C a A---.}
\milcpaso{2}{milcMagenta}{\textbf{La cuenta de los dos números (7 min).} Escribe cuántos crees que saben la versión que circula y cuántos sabrían una corrección hecha una vez. \emph{Mira la diferencia y quédate con ella:} \textbf{ese es el tamaño del daño}, y es lo que no se ve cuando uno «solo comenta».}
\milcpaso{3}{milcMagenta}{\textbf{Las tres preguntas (8 min).} Escríbelas en la contraportada del cuaderno o en una nota del celular. \textbf{Pruébalas con tres cosas} que hayas oído esta semana sobre alguien ---incluidas las que te parecieron inofensivas---. \emph{Vas a ver que casi ninguna pasa la primera pregunta.}}
\milcpaso{4}{milcMagenta}{\textbf{La corrección donde se dijo (10 min, en parejas).} Escribe la corrección de tu rumor \textbf{como si tuvieras que ponerla en el mismo lugar donde circuló}: el mismo grupo, la misma gente. \textbf{Máximo tres frases, sin regañar a nadie y sin acusar a quien lo dijo} ---si acusas, la conversación pasa a ser sobre eso y el rumor se queda intacto---. \emph{Y una advertencia honesta que hay que decir en voz alta:} \textbf{a veces esto ya no alcanza}. \emph{Cuando el daño lleva meses, la corrección no lo deshace ---y eso no es razón para no hacerla: es razón para no empezar la próxima---.}}

\begin{drawbox}{milcMagenta}{Antes de cerrar tu cadena}
\begin{checklista}
\item Rastreaste la cadena hacia atrás y escribiste \textbf{dónde se acabó}.
\item Contaste, sin adornos, \textbf{a cuántos se lo repetiste tú}.
\item Escribiste los \textbf{dos números} y miraste la diferencia.
\item Tus tres preguntas están escritas donde las vas a ver.
\item Tu corrección cabe en tres frases, va \textbf{donde fue el rumor} y no acusa a nadie.
\item Entendiste que hay daños sin remiendo, y que esa es la razón para no empezar el próximo.
\end{checklista}
\end{drawbox}

\begin{softbox}{milcMagenta}{milcGris}{Plantilla de trabajo}
\textbf{Las tres preguntas.} \emph{«¿Quién me lo dijo, y a esa persona quién? · ¿Alguien lo comprobó con ella? · Si fuera falso, ¿se podría deshacer?»} \\[1mm]
\textbf{La corrección.} \emph{«Lo que se dijo de \_\_\_\_ no era cierto. Lo comprobé con \_\_\_\_. Yo lo repetí y me equivoqué.»} \\[1mm]
\textbf{El corte.} \emph{«No sé si eso es cierto, mejor no lo repito.»} ---una frase, sin discutir con nadie---.
\end{softbox}

\begin{infoband}{milcMagenta}


\textbf{En tu cuaderno (Actividad 3 · CREA).} Título: «Actividad 3. Cortar la cadena». \textbf{Formato:} la cadena completa + los dos números + las tres preguntas probadas con tres casos + la corrección de tres frases. \textbf{Extensión:} una página. \textbf{Sabes que terminaste cuando} puedes decir \textbf{en qué punto la cadena se acaba en ti}.
\end{infoband}

\puente{El producto es una cadena rastreada. \emph{Casi nadie ha hecho este ejercicio, y casi todos se llevan la misma sorpresa: la cadena no lleva a ninguna parte.}}

\milcestacion{milcMostaza}{Producto de la sesión}
\textbf{Cadena del rumor: rastreo hacia atrás, la cuenta de los dos números y una regla de tres preguntas}.
\begin{softbox}{milcMostaza}{milcGris}{Criterios de calidad}
(1) La cadena está rastreada hacia atrás e identifica el punto exacto donde se pierde la fuente. (2) El rastreo incluye la propia participación, con el número de veces que se repitió. (3) Los dos números están calculados y se explica por qué la brecha no se cierra sola. (4) La corrección va dirigida al mismo lugar donde circuló el rumor y no acusa a quien lo dijo. (5) Se reconoce que hay daños sin remiendo, sin usarlo como excusa para no corregir.
\end{softbox}

\puente{Antes de cerrar, revisa lo incómodo: si en tu cadena tú apareces como uno de los eslabones, eso es exactamente lo que había que descubrir.}

\milcestacion{milcVino}{Evaluación liberadora · cinco dimensiones}

\begin{softbox}{milcVino}{milcGris}{Sentido de la evaluación}
Evaluar no es solo revisar una nota. Es reconocer qué aprendí, qué cambió en mí, qué emociones manejé, cómo me relaciono con la ciudad y la comunidad, y qué le dejaré a quien viene detrás.
\end{softbox}

\textbf{Mi reflexión liberadora en cinco dimensiones.} Completa cada frase iniciada:

\small
\begin{tabularx}{\textwidth}{>{\bfseries\RaggedRight\arraybackslash}p{3.4cm}Y>{\RaggedRight\arraybackslash}p{4.6cm}}
\rowcolor{milcVino}\color{white}Dimensión & \color{white}\bfseries Mi respuesta (frame starter) & \color{white}\bfseries Algo que descubrí\\
1. Personal & Lo que más cambió en mí fue \linead{6.5cm} & \linead{4.2cm}\\
2. Emocional & La emoción más fuerte fue \linead{2.5cm} y la regulé \linead{3cm} & \linead{4.2cm}\\
3. Ciudadana & En democracia esto importa porque \linead{6cm} & \linead{4.2cm}\\
4. Local (Cartago) & En mi barrio sería útil para \linead{6.5cm} & \linead{4.2cm}\\
5. Intergeneracional & A mi abuela/abuelo le diría \linead{6.5cm} & \linead{4.2cm}\\
\end{tabularx}
\normalsize

\begin{infoband}{milcVino}
\textbf{Para la próxima sustentación.} Vuelve a esta tabla en seis meses. Verás que las respuestas cambian — no porque lo aprendido cambie, sino porque tú cambias. La evaluación liberadora es una conversación contigo a través del tiempo.
\end{infoband}


\puente{Cerramos con tres voces: la persona que siempre es más que su fama, la diferencia entre la cosa y la opinión sobre la cosa, y qué pasa cuando un dato circula sin lo que lo sostenía.}

\milcestacion{milcGradeDark}{Triángulo de pensamiento · cierre filosófico}

\begin{softbox}{milcVino}{milcGris}{Dussel · lente del nosotros}
\textit{«Analéctico quiere indicar el hecho real humano por el que todo hombre, todo grupo o pueblo se sitúa siempre más allá (aná-) del horizonte de la totalidad.»}

{\footnotesize\itshape\color{milcNegro!70}--- Enrique Dussel · Filosofía de la liberación (1977), §2.4}

Una fama es \textbf{una totalidad}: un relato cerrado que pretende decir todo lo que alguien es. \emph{«Ese es el que\ldots», y ya no hace falta averiguar nada más.} \textbf{Dussel afirma justo lo contrario:} cualquier persona está \textbf{siempre más allá} de lo que el grupo alcanza a decir de ella. \emph{Y hay una prueba que se puede hacer hoy mismo:} \textbf{hablar con la persona de la que se dice algo}. \emph{Casi siempre pasa lo mismo ---y es incómodo---:} la fama no se sostiene diez minutos de conversación. \textbf{Por eso el rumor necesita que nadie hable con el aludido, y por eso preguntarle directamente es lo más eficaz que existe contra un rumor.}

\textbf{De la persona más comentada de mi colegio, ¿alguna vez he hablado con ella más de cinco minutos?}
\end{softbox}
\linead{\linewidth}\\[1mm]
\linead{\linewidth}

\begin{softbox}{milcMostaza}{milcGris}{Estoicismo · Epicteto · Enquiridión, 5 (c. 125 d.C.) · cuidado interior}
\textit{«No son las cosas las que perturban a los hombres, sino las opiniones sobre las cosas.»}

{\footnotesize\itshape\color{milcNegro!70}--- Epicteto · Enquiridión, 5 (c. 125 d.C.)}

Esta frase corta como un bisturí en los dos sentidos. \textbf{Hacia afuera:} lo que circula en un rumor \textbf{no son los hechos, son opiniones sobre los hechos} ---y entre más pasa de boca en boca, más opinión y menos hecho queda---. \textbf{Hacia adentro, y con cuidado:} \emph{no significa que a quien le pusieron una fama le duela por débil}. \textbf{Significa otra cosa, más útil:} tu valor \textbf{no lo decide el relato que circule sobre ti}. \emph{Y hay un consuelo real ahí, del que conviene acordarse a los trece años:} \textbf{la fama que hoy se siente definitiva es una opinión ---y las opiniones cambian de dueño más rápido que las personas---.}

\textbf{¿Estoy repitiendo un hecho que comprobé, o una opinión que me llegó ya empacada?}
\end{softbox}
\linead{\linewidth}\\[1mm]
\linead{\linewidth}

\begin{softbox}{milcTurquesa}{milcGris}{Floridi · lente de la infoesfera}
\textit{«Los pequeños patrones solo pueden ser significativos si se agregan correctamente, se comparan y se procesan a tiempo.»}

{\footnotesize\itshape\color{milcNegro!70}--- Luciano Floridi · Big data and their epistemological challenge (2012)}

Floridi habla de datos, pero describe exactamente cómo se fabrica una fama. \emph{Un rumor toma \textbf{pequeños patrones} ---que llegó tarde tres veces, que se sentó con alguien, que no fue a una fiesta--- y los \textbf{agrega} en una sola historia que suena redonda.} \textbf{El problema es la palabra \emph{correctamente}:} nadie comparó esos datos con nada, nadie preguntó, nadie miró lo que faltaba. \emph{Y en un chat es peor, porque el mensaje corto llega sin lo que lo sostenía} ---igual que el titular sin el artículo---. \textbf{Regla práctica:} \emph{si una historia sobre alguien te suena demasiado redonda, probablemente le quitaron lo que no cuadraba.}

\textbf{La última historia que oí sobre alguien, ¿qué datos habría necesitado para saber si era cierta?}
\end{softbox}
\linead{\linewidth}\\[1mm]
\linead{\linewidth}

\begin{infoband}{milcNegro}
\textbf{Nota para el docente.} Las tres son citas textuales y llevan su referencia debajo. Puedes remitir a la obra en clase.
\end{infoband}

\milcestacion{milcVino}{Mi autoevaluación · marca una casilla por fila}

{\small
\setlength{\extrarowheight}{2.2pt}
\begin{tabularx}{\textwidth}{>{\RaggedRight\arraybackslash\bfseries}p{3.0cm}YYY}
\rowcolor{milcVino}\color{white}\bfseries Criterio & \color{white}\bfseries Lo logré & \color{white}\bfseries En proceso & \color{white}\bfseries Necesito apoyo\\[.6mm]
Comprensión del tema & \milcopcion{Explico las ideas con mis palabras.} & \milcopcion{Reconozco las ideas, falta precisión.} & \milcopcion{Necesito repasar lo básico.}\\[.8mm]
\rowcolor{milcGris}Calidad del rastreo & \milcopcion{Reconstruí la cadena hacia atrás y entendí por qué la corrección nunca alcanza al rumor.} & \milcopcion{Sé que el rumor hace daño, pero creo que basta con no inventarlo yo.} & \milcopcion{Sigo pensando que repetir algo que ya sabe todo el mundo no cuenta.}\\[.8mm]
Aplicación y producto & \milcopcion{Mi producto cumple los criterios.} & \milcopcion{Está armado, con ajustes pendientes.} & \milcopcion{Aún está en construcción.}\\[.8mm]
\rowcolor{milcGris}Reflexión 5 dimensiones & \milcopcion{Respondí cada una con honestidad.} & \milcopcion{Respondí algunas con detalle.} & \milcopcion{Mis respuestas son muy generales.}\\[.8mm]
Triángulo de pensamiento & \milcopcion{Conecté las 3 voces con mi trabajo.} & \milcopcion{Conecté al menos una.} & \milcopcion{No las relaciono con mi proceso.}\\[.8mm]
\end{tabularx}}

\vspace{3mm}
\textbf{Hoy entendí que:}\\[1mm]
\linead{\linewidth}\\[1mm]
\linead{\linewidth}

\textbf{Mi compromiso para la próxima sesión es:}\\[1mm]
\linead{\linewidth}\\[1mm]
\linead{\linewidth}

\begin{softbox}{milcMostaza}{milcGris}{Cita de despedida · Marco Aurelio}
\textit{``El obstáculo en el camino se vuelve el camino. Lo que estorba al camino se vuelve camino.''}\\[1mm]
Cada error de hoy, bien observado, se convierte en pieza del camino que pisarás mañana.
\end{softbox}

\small
\begin{tabularx}{\textwidth}{>{\bfseries}p{3.6cm}Y}
\rowcolor{milcGradeDark!12}Nombre completo: & \linead{10cm}\\
Fecha: & \linead{4cm} \quad Curso/grupo: \linead{4cm}\\
Mi firma: & \linead{9cm}\\
\end{tabularx}
\normalsize

\begin{infoband}{milcGradeDark}
\textbf{Para la próxima sesión.} Dos cosas. \textbf{Primero, corta una cadena} ---una sola vez: cuando alguien vaya a contarte algo sobre otro, di tu frase y no lo repitas---. \emph{Anota qué se sintió; casi todo el mundo reporta lo mismo: incomodidad de dos segundos y después nada.} \textbf{Y si te alcanza el valor, habla con la persona de la que se dice algo} ---no del rumor: de cualquier cosa---, \textbf{y anota qué aprendiste de ella que no cuadraba con su fama}. \textbf{Segundo, pregunta en tu casa cómo se cuidaba el nombre:} averigua con alguien mayor \textbf{qué significaba «tener buen nombre» en su época} ---qué se podía perder por un chisme, si a una familia le cerraban puertas por hablas, cómo se defendía alguien de una fama---. \textbf{Trae:} qué se sintió cortar la cadena y lo que te contaron. \emph{Vas a encontrar una idea antigua y muy exacta: el nombre era algo que se cuidaba como se cuida una herramienta de trabajo, porque de él dependían cosas concretas ---el fiado, el trabajo, el matrimonio---.}
\end{infoband}

\end{document}
